\documentclass{beamer}
\usepackage[utf8]{inputenc}
\usepackage[spanish]{babel}
\usepackage{listings}
\usepackage{xcolor}
\usepackage{tikz}
\usepackage{tcolorbox}
\usepackage{fontawesome5}

% Configuración de colores - Gruvbox Light con marrón principal
\definecolor{bg0}{HTML}{fbf1c7}
\definecolor{bg1}{HTML}{ebdbb2}
\definecolor{bg2}{HTML}{d5c4a1}
\definecolor{bg3}{HTML}{bdae93}
\definecolor{fg0}{HTML}{282828}
\definecolor{fg1}{HTML}{3c3836}
\definecolor{fg2}{HTML}{504945}
\definecolor{aqua}{HTML}{689d6a}
\definecolor{aquabright}{HTML}{427b58}
\definecolor{blue}{HTML}{458588}
\definecolor{red}{HTML}{cc241d}
\definecolor{orange}{HTML}{d65d0e}
\definecolor{yellow}{HTML}{d79921}
\definecolor{green}{HTML}{98971a}
\definecolor{purple}{HTML}{b16286}
\definecolor{brown}{HTML}{AB886D}
\definecolor{brownbright}{HTML}{8B6849}

% Colores para código
\definecolor{codegreen}{HTML}{98971a}
\definecolor{codegray}{HTML}{7c6f64}
\definecolor{codepurple}{HTML}{b16286}
\definecolor{backcolour}{HTML}{f9f5d7}

% Configuración de código
\lstdefinestyle{pythonstyle}{
    backgroundcolor=\color{backcolour},   
    commentstyle=\color{codegreen},
    keywordstyle=\color{brownbright}\bfseries,
    numberstyle=\tiny\color{codegray},
    stringstyle=\color{codepurple},
    basicstyle=\ttfamily\footnotesize\color{fg0},
    breakatwhitespace=false,         
    breaklines=true,                 
    captionpos=b,                    
    keepspaces=true,                 
    numbers=left,                    
    numbersep=3pt,                  
    showspaces=false,                
    showstringspaces=false,
    showtabs=false,                  
    tabsize=2,
    language=Python,
    frame=single,
    rulecolor=\color{brown},
    xleftmargin=8pt,
    xrightmargin=8pt,
    inputencoding=utf8,
    extendedchars=true,
    literate=
        {á}{{\'a}}1 {é}{{\'e}}1 {í}{{\'i}}1 {ó}{{\'o}}1 {ú}{{\'u}}1
        {Á}{{\'A}}1 {É}{{\'E}}1 {Í}{{\'I}}1 {Ó}{{\'O}}1 {Ú}{{\'U}}1
        {ñ}{{\~n}}1 {Ñ}{{\~N}}1
        {¿}{{?`}}1 {¡}{{!`}}1
}

\lstset{style=pythonstyle}

% Tema personalizado con marrón principal
\usetheme{Madrid}
\setbeamercolor{structure}{fg=brown}
\setbeamercolor{palette primary}{bg=brown,fg=bg0}
\setbeamercolor{palette secondary}{bg=brownbright,fg=bg0}
\setbeamercolor{palette tertiary}{bg=fg1,fg=bg0}
\setbeamercolor{palette quaternary}{bg=brown,fg=bg0}
\setbeamercolor{block title}{bg=brown,fg=bg0}
\setbeamercolor{block body}{bg=bg1,fg=fg0}
\setbeamercolor{frametitle}{bg=brown,fg=bg0}
\setbeamercolor{title}{fg=fg0}
\setbeamercolor{subtitle}{fg=fg1}
\setbeamercolor{normal text}{fg=fg0,bg=bg0}
\setbeamercolor{item}{fg=brown}
\setbeamercolor{section in toc}{fg=brown}

% Cajas personalizadas con marrón
\newtcolorbox{conceptbox}[1]{
    colback=bg1,
    colframe=brown,
    fonttitle=\bfseries\small,
    title=#1,
    coltitle=fg0,
    left=3pt,
    right=3pt,
    top=3pt,
    bottom=3pt,
    boxsep=2pt
}

\newtcolorbox{notebox}{
    colback=yellow!15,
    colframe=orange,
    leftrule=2mm,
    rightrule=0mm,
    toprule=0mm,
    bottomrule=0mm,
    arc=0mm,
    left=3pt,
    right=3pt,
    top=2pt,
    bottom=2pt
}

\title{\textbf{Introducción a la Programación con Python}}
\subtitle{Semana 1: Condicionales y Control de Flujo}
\author{$\nabla$Nablabs}
\date{}
\begin{document}

\begin{frame}
\titlepage
\end{frame}

\begin{frame}
\frametitle{Contenido}
\tableofcontents
\end{frame}

\section{Operadores de comparación}

\begin{frame}[fragile]{Operadores de comparación \ \faBalanceScale}
\begin{conceptbox}{¿Para qué sirven?}
Comparan valores y devuelven un resultado booleano (True/False)
\end{conceptbox}

\begin{lstlisting}
5 > 3      # True
2 < 1      # False
4 >= 4     # True
7 <= 10    # True
3 == 3     # True
2 != 5     # True
\end{lstlisting}

\vspace{-0.2cm}
\begin{tcolorbox}[colback=green!10,colframe=green,title={\small\faArrowCircleRight \ Salida esperada},coltitle=fg0,left=3pt,right=3pt,top=2pt,bottom=2pt,boxsep=1pt]
\small\texttt{True\\
False\\
True\\
True\\
True\\
True}
\end{tcolorbox}

\vspace{-0.1cm}
\begin{notebox}
\small\faInfoCircle \ El resultado de una comparación es siempre True o False
\end{notebox}
\end{frame}

\section{Sentencias if, elif, else}

\begin{frame}[fragile]{Estructura if, elif, else \ \faCodeBranch}
\begin{conceptbox}{¿Para qué sirven?}
Permiten ejecutar código solo si se cumple una condición
\end{conceptbox}

\begin{lstlisting}
edad = 18
if edad >= 18:
    print("Eres mayor de edad")
elif edad >= 13:
    print("Eres adolescente")
else:
    print("Eres niño/a")
\end{lstlisting}

\vspace{-0.2cm}
\begin{tcolorbox}[colback=green!10,colframe=green,title={\small\faArrowCircleRight \ Salida},coltitle=fg0,left=3pt,right=3pt,top=2pt,bottom=2pt,boxsep=1pt]
\small\texttt{Eres mayor de edad}
\end{tcolorbox}

\vspace{-0.1cm}
\begin{notebox}
\small\faExclamationTriangle \ El bloque else es opcional
\end{notebox}
\end{frame}

\section{Operadores lógicos}

\begin{frame}[fragile]{Operadores lógicos \ \faProjectDiagram}
\begin{conceptbox}{¿Para qué sirven?}
Combinan condiciones para tomar decisiones más complejas
\end{conceptbox}

\begin{lstlisting}
edad = 20
es_estudiante = True

if edad >= 18 and es_estudiante:
    print("Descuento para estudiantes adultos")

if edad < 18 or es_estudiante:
    print("Descuento disponible")

if not es_estudiante:
    print("No eres estudiante")
\end{lstlisting}

\vspace{-0.2cm}
\begin{tcolorbox}[colback=green!10,colframe=green,title={\small\faArrowCircleRight \ Salida},coltitle=fg0,left=3pt,right=3pt,top=2pt,bottom=2pt,boxsep=1pt]
\small\texttt{Descuento para estudiantes adultos\\
Descuento disponible}
\end{tcolorbox}
\end{frame}

\section{Valores y expresiones booleanas}

\begin{frame}[fragile]{Valores booleanos \ \faToggleOn}
\begin{conceptbox}{True y False}
Son los dos únicos valores booleanos en Python
\end{conceptbox}

\begin{lstlisting}
mayor = 10 > 5      # True
igual = 3 == 7      # False
activo = True
inactivo = False
\end{lstlisting}

\vspace{-0.2cm}
\begin{tcolorbox}[colback=green!10,colframe=green,title={\small\faArrowCircleRight \ Salida esperada},coltitle=fg0,left=3pt,right=3pt,top=2pt,bottom=2pt,boxsep=1pt]
\small\texttt{mayor = True\\
igual = False\\
activo = True\\
inactivo = False}
\end{tcolorbox}

\vspace{-0.1cm}
\begin{notebox}
\small\faLightbulb \ Las expresiones de comparación devuelven booleanos
\end{notebox}
\end{frame}

\section{Anidamiento de condicionales}

\begin{frame}[fragile]{Condicionales anidados \ \faSitemap}
\begin{conceptbox}{¿Para qué anidar?}
Para tomar decisiones dentro de otras decisiones
\end{conceptbox}

\begin{lstlisting}
edad = 16
if edad >= 13:
    if edad >= 18:
        print("Adulto")
    else:
        print("Adolescente")
else:
    print("Niño/a")
\end{lstlisting}

\vspace{-0.2cm}
\begin{tcolorbox}[colback=green!10,colframe=green,title={\small\faArrowCircleRight \ Salida},coltitle=fg0,left=3pt,right=3pt,top=2pt,bottom=2pt,boxsep=1pt]
\small\texttt{Adolescente}
\end{tcolorbox}
\end{frame}

\section{Operador match (switch)}

\begin{frame}[fragile]{Operador match \ \faRandom}
\begin{conceptbox}{¿Qué es match?}
Permite comparar un valor contra varios casos (Python 3.10+)
\end{conceptbox}

\begin{lstlisting}
comando = "saludar"
match comando:
    case "saludar":
        print("¡Hola!")
    case "despedir":
        print("Adiós!")
    case _:
        print("Comando no reconocido")
\end{lstlisting}

\vspace{-0.2cm}
\begin{tcolorbox}[colback=green!10,colframe=green,title={\small\faArrowCircleRight \ Salida},coltitle=fg0,left=3pt,right=3pt,top=2pt,bottom=2pt,boxsep=1pt]
\small\texttt{¡Hola!}
\end{tcolorbox}

\vspace{-0.1cm}
\begin{notebox}
\small\faInfoCircle \ El operador match es similar a switch en otros lenguajes
\end{notebox}
\end{frame}

\section{Ejemplo integrador}

\begin{frame}[fragile]{Ejemplo práctico \ \faUserCheck}
\begin{lstlisting}
edad = int(input("Edad: "))
estudiante = input("¿Eres estudiante? (s/n): ").strip().lower() == "s"

if edad >= 18:
    if estudiante:
        print("Eres un estudiante adulto")
    else:
        print("Eres un adulto no estudiante")
elif edad >= 13:
    print("Eres adolescente")
else:
    print("Eres niño/a")
\end{lstlisting}

\vspace{0.2cm}
\begin{columns}[T]
\column{0.48\textwidth}
\begin{tcolorbox}[colback=blue!10,colframe=blue,title={\small\faArrowDown \ Entrada},coltitle=fg0,left=2pt,right=2pt,top=2pt,bottom=2pt]
\small\texttt{Edad: 20\\
¿Eres estudiante? (s/n): s}
\end{tcolorbox}

\column{0.48\textwidth}
\begin{tcolorbox}[colback=green!10,colframe=green,title={\small\faArrowUp \ Salida},coltitle=fg0,left=2pt,right=2pt,top=2pt,bottom=2pt]
\small\texttt{Eres un estudiante adulto}
\end{tcolorbox}
\end{columns}
\end{frame}

\begin{frame}{Resumen \ \faListUl}
\begin{conceptbox}{¿Qué aprendimos hoy?}
\small
\begin{itemize}
    \item[\faCheck] Operadores de comparación
    \item[\faCheck] Sentencias if, elif, else
    \item[\faCheck] Operadores lógicos
    \item[\faCheck] Valores y expresiones booleanas
    \item[\faCheck] Anidamiento de condicionales
    \item[\faCheck] Operador match
\end{itemize}
\end{conceptbox}

\vspace{0.3cm}
\begin{center}
\large\textcolor{brown}{\faCode \ ¡Practica con condicionales!}
\end{center}
\end{frame}

\begin{frame}
\begin{center}
\Huge\textcolor{brown}{\faQuestionCircle}\\
\vspace{0.5cm}
\Huge ¡Preguntas!\\
\vspace{1cm}
\Large\textcolor{brownbright}{Gracias por su atención}
\end{center}
\end{frame}

\end{document}
